%% -*- coding: utf-8 -*-
\documentclass[12pt,pagesize,paper=192mm:108mm]{scrbook} 
%1920x1080 1280x720
\areaset[current]{192mm}{108mm}
\usepackage{calc}
\usepackage[T2A]{fontenc}
\usepackage[utf8]{inputenc}
\usepackage[english,russian]{babel}
\usepackage{microtype}
\usepackage{misccorr}
\usepackage{cmap}
%\usepackage[unicode=true]{hyperref}
\usepackage{graphicx}
\usepackage{amssymb}
\usepackage{amsmath}
%\usepackage{srcltx}
\usepackage{textcomp}
\usepackage{xspace}
%научные символы и смайлики \smiley \frownie
\usepackage{wasysym}
\usepackage{ccicons}
\begin{document}
\begin{titlepage}
  \vspace*{-0.5em}
  \begin{center}    
    \hspace*{3em}
    \begin{minipage}[t]{3em}
      \includegraphics[width=\textwidth]{../BINP-logo}
    \end{minipage}\hfill
    \begin{minipage}{0.23\linewidth}
    \includegraphics[width=\textwidth]{../NSU-logo}
    \end{minipage}
    \hfill
    \hspace*{6em}

    Кафедра теоретической физики физического факультета НГУ
    \medskip

    \Large
    Профессор Черняк В.\,Л.
    \bigskip

    \huge
    \textbf{Теория электрослабых взаимодействий}
    \bigskip

    \Large
    Лекция № 11
    \vfill

    % \normalsize
    % \begin{minipage}{0.65\linewidth}
    % \end{minipage}
    \vfill

\normalsize    Новосибирск 2013
  \smallskip

  \ccbysa
  \end{center}
\end{titlepage}
\newpage

\vspace*{-1em}
\begin{center}
 \vfill
  \begin{minipage}{0.66\linewidth}
    Нахождение фундаментальных параметров сильной и электрослабой
    теории. Полулептонные процессы. Бета"=распад $\pi$"=мезона: ширина
    распада $\pi^-\to\pi^0 e^-\bar{\nu}_e$.  Вычисление адронного
    матричного элемента, входящего в амплитуду этого распада.
    Формфакторы $f_+(q^2)$ и $f_-(q^2)$.  Изотопическая симметрия и
    дивергенции векторных токов. Подавление формфактора $f_-$ в
    изотопическом пределе.  Вычисление формфактора
    $f_+(q^2=0)=\sqrt{2}$ в изотопическом пределе.  Анализ ответа для
    ширины распада $\pi^-\to\pi^0 e^-\bar{\nu}_e$: подавление
    множителем из фазового объёма, независимое измерение $V_{ud}$,
    пион"=бета эксперимент.  Распад $K^-\to\pi^0 e^-\bar{\nu}_e$:
    точность $SU(3)$ флейворной симметрии при вычислении $f_+(q^2=0)$.
  \end{minipage}
  \vfill
  % Новосибирск 2013

  % \normalsize \ccbysa
\end{center}
\end{document}
